\documentclass[../../../../main.tex]{subfiles}
\begin{document}

\begin{flushright}
\footnotesize\textit{【自動文字起こし・要確認】}
\end{flushright}

[解]\\
(1) $g(x)-x = f(x)^2+af(x)+b-x$\\
$\quad= \{f(x)-x\}\{f(x)+a+x+1\} \quad \qed$

(2) $g(p)=p \land f(p) \neq p$ の時（$p \in \mathbb{R}$）(1)から
\[
\begin{cases}
f(p)+a+p+1=0 \\
f(p) \neq p
\end{cases}
\]
\[
\Leftrightarrow
\begin{cases}
p^2+(a+1)p+a+b+1=0 & \cdots \text{①} \\
p^2+(a-1)p+b \neq 0 & \cdots \text{②}
\end{cases}
\]
①をみたす$p$がある条件は判別式$D$として
\[ D = (a+1)^2-4(a+b+1) \ge 0 \quad \cdots \text{③} \]
このもとで、「①の左辺と②の左辺が2つとも同じ解を持つ（重解含む）」がみたされなければよい。

$1^\circ \ D>0$ の時\\
①は2実数解を持ち、②も2次方程式だから①②が一致しなければよく、係数比較して
\[ a+1 \neq a-1, \quad a+b+1 \neq b \]
前者が常にみたされるから$D>0$は適。

$2^\circ \ D=0$ の時\\
①は重解$p=-\frac{a+1}{2}$を持つから、これが②をみたさなければよい。代入して
\[ -\frac{a+1}{2}\left(-\frac{a+1}{2}+a-1\right)+\frac{1}{4}a^2-\frac{1}{2}a-\frac{3}{4}=0 \]
これは$a$についての恒等式だから$D=0$の時、これが常にみたされ、不適。

以上から、条件は$D>0$で、図示して下図斜線部（境界含まず）

\begin{center}
\begin{tikzpicture}[scale=0.8]
    \fill[pattern=north east lines, pattern color=gray!50] (-3, -3) -- plot[domain=-3:4] (\x, {0.25*\x*\x - 0.5*\x - 0.75}) -- (4, -3) -- cycle;
    \draw[->] (-3.5, 0) -- (4.5, 0) node[right] {$a$};
    \draw[->] (0, -3.5) -- (0, 3.5) node[above] {$b$};
    \draw[domain=-3.5:4.5, smooth, variable=\x] plot ({\x}, {0.25*\x*\x - 0.5*\x - 0.75}) node[right] {$b=\frac{1}{4}a^2-\frac{1}{2}a-\frac{3}{4}$};
    \draw[dashed] (1, 0) -- (1, -1) -- (0, -1);
    \node[below] at (1, 0) {$1$};
    \node[left] at (0, -1) {$-1$};
    \node[above right] at (0, 0) {$O$};
\end{tikzpicture}
\end{center}

\end{document}
